\DocumentMetadata{
  pdfversion=2.0,
  pdfstandard=ua-2,
  lang=en-US,
  tagging=on,
  tagging-setup={math/setup=mathml-SE,tabsorder=structure}
}
\documentclass[11pt,letterpaper]{article}
\usepackage[margin=0.68in,includefoot]{geometry}
\usepackage{newcomputermodern}
\usepackage{amsmath,amscd,mathtools}
\providecommand{\square}{\mdlgwhtsquare}
\usepackage[hidelinks]{hyperref}
\hypersetup{
  pdftitle={Analysis First-Year Exam, January 2024, Part 2},
  pdfauthor={University of Florida Department of Mathematics},
  pdfsubject={Graduate mathematics examination},
  pdfdisplaydoctitle=true,
  bookmarksopen=true
}
\setcounter{secnumdepth}{0}
\setlength{\parindent}{0pt}
\setlength{\parskip}{0.28em}
\setlength{\emergencystretch}{3em}
\setlength{\leftmargini}{2.15em}
\setlength{\leftmarginii}{2.65em}
\setlength{\leftmarginiii}{3em}
\renewcommand{\labelenumi}{\textbf{\theenumi.}}
\pagestyle{empty}
\raggedbottom
\allowdisplaybreaks
\newenvironment{examproblems}[1]{%
  \begin{enumerate}%
  \setcounter{enumi}{#1}%
  \setlength{\topsep}{0.35em}%
  \setlength{\partopsep}{0pt}%
  \setlength{\itemsep}{0.48em}%
  \setlength{\parsep}{0.18em}%
}{\end{enumerate}}
\newenvironment{examsubparts}{%
  \begin{enumerate}%
  \setlength{\topsep}{0.18em}%
  \setlength{\partopsep}{0pt}%
  \setlength{\itemsep}{0.16em}%
  \setlength{\parsep}{0.1em}%
}{\end{enumerate}}
\newenvironment{examinstructions}{%
  \par\begingroup\itshape
}{%
  \par\endgroup\vspace{0.18em}
}
\makeatletter
\renewcommand\section{\@startsection{section}{1}{0pt}%
  {-1.25ex plus -.25ex minus -.1ex}%
  {0.55ex plus .1ex}%
  {\normalfont\large\bfseries}}
\renewcommand{\@maketitle}{%
  \newpage\null\vskip -1.2em%
  {\footnotesize\bfseries
    \href{https://www.ufl.edu/}{University of Florida}, \href{https://math.ufl.edu/}{Department of Mathematics}\par}%
  \vskip 0.18em%
  \tagstructbegin{tag=Title}%
    {\LARGE\bfseries\href{https://gma.math.ufl.edu/exams/analysis-first-year/}{Analysis First-Year Exam}, January 2024, Part 2\par}%
  \tagstructend%
  \vskip 0.35em%
  \tagmcbegin{artifact}%
    \hrule height 1pt%
  \tagmcend%
  \vskip 0.55em%
}
\makeatother
\title{Analysis First-Year Exam, January 2024, Part 2}
\author{}
\date{}
\begin{document}
\maketitle
\thispagestyle{empty}
\begin{examinstructions}
Answer FOUR questions in detail. State carefully any results used without proof.
\end{examinstructions}
\begin{examproblems}{0}
\item
Prove that the function \(f:[a,b]\to\mathbb{R}\) is Riemann integrable iff for each \(\varepsilon>0\) there exist Riemann integrable functions \(\ell,u:[a,b]\to\mathbb{R}\) such that \(\ell\leq f\leq u\) pointwise on \([a,b]\) and \(\int_a^b(u-\ell)<\varepsilon\).
\item
For each positive integer~\(n\) let \(f_n:[0,1]\to[0,1]\) be a Riemann integrable function. Which (if any) of the following conditions is sufficient to ensure that \(\int_0^1 f_n(t)\,dt\to0\) as \(n\to\infty\)?
\begin{examsubparts}
\item[\textnormal{(i)}]
\(f_n\to0\) uniformly on \([0,1]\);
\item[\textnormal{(ii)}]
\(f_n\to0\) pointwise on \([0,1]\).
\end{examsubparts}
\item
To each continuous function \(f:[0,1]\to\mathbb{R}\) associate the sequence \((a(f)_n)_{n=0}^{\infty}\) defined by 
\[
a(f)_n=\int_0^1 f(t)t^n\,dt.
\]
 Prove that the map \(f\mapsto a(f)\) is injective.
\item
Let the bounded set \(X\subseteq\mathbb{R}\) be Lebesgue measurable. Prove that:
\begin{examsubparts}
\item[\textnormal{(i)}]
for each \(\varepsilon>0\) there exists an open set \(U\supseteq X\) whose Lebesgue measure \(\lambda(U)\) is less than \(\lambda(X)+\varepsilon\);
\item[\textnormal{(ii)}]
there exists a Borel set \(B\supseteq X\) such that \(\lambda(B)=\lambda(X)\).
\end{examsubparts}
\item
Let \((f_n)_{n=0}^{\infty}\) be a sequence of functions on the measurable space \((X,\mathcal{A})\) with values in \([0,\infty)\). Prove that each of the following subsets of~\(X\) is measurable:
\begin{examsubparts}
\item[\textnormal{(i)}]
\(B=\{x:\text{ the sequence }(f_n(x))_{n=0}^{\infty}\text{ is bounded}\}\);
\item[\textnormal{(ii)}]
\(C=\{x:\text{ the series }\sum_{n=0}^{\infty}f_n(x)\text{ is convergent}\}\).
\end{examsubparts}
\end{examproblems}
\end{document}
